% GROK REVIEW::GROK-0512::2026-10-14T05:12:00+09:00
\documentclass[11pt]{article}
\usepackage[margin=1in]{geometry}
\usepackage{amsmath,amssymb}
\title{Independent Grok Review --- 2026-10-14 05:12 KST}
\author{Grok}
\date{14 October 2026}

\begin{document}
\maketitle

\begin{center}
\fbox{\parbox{0.9\linewidth}{
\textbf{Executive status: P versus NP remains OPEN.}
No load-bearing implication is closed by this review.
}}
\end{center}

\section{Scope}
This independent review examines the living document at repository-root \texttt{main.tex} (ROOTGLUE1 synthesis, reconciled 2026-10-14), the specialist reports TRACEISO1 and related affine-residual materials, and the new candidate interfaces for affine root-list normal forms and three-root joint compatibility. It confirms that all current firewalls remain correctly classified and that the unrestricted whole-computation lower bound stays OPEN.

\section{Confirmation of current spine}
The affine base, residual contraction, cotree witness bijection, and upper-bound escapes remain PROVED/EXACT. Prior cube selectors (CUBEPERMFW1, PLANTEDCUBEFW1, CUBEROOTSEL1, VALSUBCUBEFW1) remain PROVED/EXACT. TRACEISO1's deterministic root-isolation theorem for every explicit binary-field presentation is confirmed as EXACT by independent ROOTGLUE1 synthesis review; it is consistent with known small-characteristic factoring algorithms and does not claim hidden-field recovery from arbitrary ANF. The new affine root compatibility normal form (including sets of size at most two) and the F4 three-root encoding of graph 3-colouring remain CANDIDATE pending independent specialist review. The logarithmic-subfield and CYCLOROOTFW1 selectors remain valid specialized algorithms, not hardness boundaries. Growing-factor amplification remains CONDITIONAL. The unrestricted whole-computation lower bound remains OPEN.

No contradiction with prior exact results is found. The three-root construction correctly relocates known NP-complete compatibility onto an explicit affine carrier of treewidth two; it is not hardness amplification. Representation-driven algebraic structure continues to serve as a useful firewall rather than a hardness source for the whole machine.

\section{Open load-bearing arrows}
The implication from an arbitrary eligible short program (or uniform decider) to superpolynomial cost on a structured residual object remains OPEN. No model-matched all-algorithm theorem exists. The firewalls relativize and algebrize; no nonrelativizing hinge is known. Opaque field representations recovered from raw quadratic ANF, arbitrary joint nonlinear compatibility, and the historical PC22 compiler remain outside the current selector theorems without additional bridges.

\section{Highest-value next target}
The single highest-value next assignment remains appropriate: build and independently test one correlation-preserving elimination rule for the explicit three-root compatibility family, measuring actual representation growth and full witness reconstruction. The rule must preserve polynomial solving of affine root domains and reject $K_4$ despite nonempty local domains. Failure of that one rule is not a lower bound for other algorithms. A successful selector or elimination is a useful firewall; it does not close a whole-machine implication.

\section{Status and recommendations}
P versus NP remains open. The living document is updated to record this independent review. All claims continue to be classified under the allowed statuses. The comments folder receives this review artifact.

% END_OF_VERIFIED_REVIEW::GROK::GROK-0512
\end{document}
